\chapter{Geometric Class Field Theory}\label{section:gcft}

The purpose of this chapter is to deduce the reduced form of geometric class
field theory, \Cref{thm:GCFT_reduced}, from the ramification calculation
established in Chapter~2.

Let $C$ be a smooth, projective, geometrically connected curve over the
perfect field $k$, let $\mdls$ be a modulus on $C$, and put
$U=C\setminus\mdls$. Given a rank-one local system $\cF$ on $U$ whose
ramification is bounded by $\mdls$, our goal is to descend its symmetric
power $\cF^{[d]}$ along the Abel--Jacobi morphism
\[
    \Phi_d\colon U^{(d)}\longrightarrow\PicCm[d]
\]
for all sufficiently large $d$.

When $\mdls=0$, the fibers of the Abel--Jacobi morphism are projective
spaces. Their trivial \'etale fundamental groups allow the desired descent
directly. When $\mdls>0$, however, the fibers of $\Phi_d$ are affine spaces,
so this argument cannot be applied to $\Phi_d$ itself. Following Takeuchi,
we resolve this difficulty by extending $\Phi_d$ to a projective space
bundle
\[
    \widetilde\Phi_d\colon
    \widetilde C_{\mdls}^{(d)}
    \longrightarrow \PicCm[d],
\]
where $\widetilde C_{\mdls}^{(d)}$ is an open subscheme of the blowup of
$C^{(d)}$ along the locus obtained by adding the modulus.

The proof then has three remaining steps. First, we use the single-point
ramification calculation of Chapter~2, together with compatibility under
products and addition of divisors, to prove that $\cF^{[d]}$ is tamely
ramified along the boundary
\[
    H=\widetilde C_{\mdls}^{(d)}\setminus U^{(d)}.
\]
Second, on the geometric generic Abel--Jacobi fiber, the boundary is the
hyperplane at infinity. A Tendler-style prime-to-$p$ power argument,
together with the structure of the abelian fundamental group of affine
space, forces the corresponding boundary inertia representation to be
trivial. Purity then extends $\cF^{[d]}$ uniquely across the boundary.
Finally, since $\widetilde\Phi_d$ is a projective space bundle, it induces
an isomorphism on \'etale fundamental groups. The extended local system
therefore descends uniquely to $\PicCm[d]$, proving
\Cref{thm:GCFT_reduced}.

We begin by recalling the generalized Picard scheme and the geometry of the
Abel--Jacobi morphism. We then introduce Takeuchi's compactification, prove
the required ramification and extension results, and conclude with the
descent argument.

\section{Generalized Picard Schemes and the Abel--Jacobi Map}\label{section:generalized-picard-and-abel-jacobi}

\subsection{The generalized Picard scheme}
We recall the notion of generalized Jacobian varieties and study their fundamental properties. The material presented here is primarily adapted from \cite{Guignard2018} and \cite{takeuchi2019blow}. For further background on the general theory of abelian varieties and Jacobians, the reader may also consult \cite{milneAV}.
Although the main theorem is over the fixed perfect field $k$, we recall the
generalized Picard construction temporarily over an arbitrary base scheme
$S$.
Let $S$ be a scheme and let $p:C\to S$ be a projective smooth morphism whose geometric fibers are connected and of dimension 1. Let $\mdls$ be a modulus on $C$, defined as an effective Cartier divisor of $C/S$ 
(i.e., a closed subscheme of $C$ which is finite flat of finite presentation over $S$).
We denote the projection $C \times_S T \to T$ by $\text{pr}$ for any $S$-scheme $T$.

\subsection*{The Functor of Points}
Let $d$ be an integer. For an $S$-scheme $T$, we consider the set of data $(\cL, \psi)$ where:
\begin{itemize}
    \item $\cL$ is an invertible sheaf of degree $d$ on $C_T$.
    \item $\psi: \Oo_{\mdls_T} \xrightarrow{\sim} \cL|_{\mdls_T}$ is a trivialization of $\cL$ along the modulus.
\end{itemize}
Two such pairs $(\cL, \psi)$ and $(\cL', \psi')$ are said to be isomorphic if there exists an isomorphism of invertible sheaves $f : \cL \to \cL'$ such that the following diagram commutes:
\[
\begin{tikzcd}
& \Oo_{\mdls_T} \arrow[dl, "\psi"'] \arrow[dr, "\psi'"] & \\
\cL|_{\mdls_T} \arrow[rr, "f|_{\mdls_T}"'] & & \cL'|_{\mdls_T}
\end{tikzcd}
\]
We define the presheaf $\text{Pic}_{C, \mdls}^{d, \text{pre}}$ on $\text{Sch}/S$ by assigning to $T$ the set of isomorphism classes of such pairs. Let $\PicCm[d]$ denote the \'etale sheafification of this presheaf.

\subsection*{Representability and Structure}
The fundamental properties of this functor are as follows:
\begin{enumerate}
    \item $\text{Pic}_{C, \mdls}^{d}$ is represented by an $S$-scheme. (Note: If $\mdls$ is faithfully flat over $S$, the presheaf is already a \'etale sheaf).
    \item $\text{Pic}_{C, \mdls}^{0}$ is a smooth commutative group $S$-scheme with geometrically connected fibers, referred to as the \textit{generalized Jacobian variety} of $C$ with modulus $\mdls$.
    \item For any $d$, $\text{Pic}_{C, \mdls}^{d}$ is a $\text{Pic}_{C, \mdls}^{0}$-torsor.
\end{enumerate}
In the case where $\mdls = 0$, we recover the standard Jacobian variety, denoted simply as $\text{Pic}_C^d$.
\subsection*{Relation to the Standard Jacobian}
We now examine the behavior of the generalized Picard scheme under the variation of the modulus. By viewing the structure along the modulus as an additional rigidification, we obtain natural transition maps corresponding to the inclusion of moduli.

Let $\mdls_1$ and $\mdls_2$ be moduli such that $\mdls_1 \subset \mdls_2$. There exists a natural map 
\[
\text{Pic}_{C, \mdls_2}^d \to \text{Pic}_{C, \mdls_1}^d
\]
obtained by restricting the rigidification from $\mdls_2$ to $\mdls_1$.
More precisely, it is an epimorphism for the \'etale topology: if
$(\cL,\psi_1)$ is an object over an $S$-scheme $T$, then, after an \'etale
cover $T'\to T$, the rigidification $\psi_1$ of
$\cL|_{(\mdls_1)_{T'}}$ extends to a rigidification of
$\cL|_{(\mdls_2)_{T'}}$. Thus every local section of
$\text{Pic}_{C,\mdls_1}^d$ lifts \'etale-locally to a section of
$\text{Pic}_{C,\mdls_2}^d$.
In particular, for any modulus $\mdls$, there is a natural surjective morphism of \'etale sheaves:
\[
\text{Pic}_{C, \mdls}^d \to \text{Pic}_C^d.
\]


\subsection*{Local Freeness and Base Change}
Let $\mdls$ be a modulus which is everywhere strictly positive. Let $g$ denote the genus of $C$, which is a locally constant function on $S$. We restrict our attention to degrees $d$ satisfying the condition:
\begin{equation} \label{equation:degree-condition}
    d \geq \max\{2g - 1 + \deg \mdls, \deg \mdls\}.
\end{equation}

Assuming $S$ is quasi-compact, such a $d$ always exists.

Fix an integer $d$ satisfying the condition above. Let $T$ be an $S$-scheme and let $\cL$ be an invertible sheaf of degree $d$ on $C_T$. One can show that the pushforwards $\mathrm{pr}_* \cL$ and $\mathrm{pr}_* \cL(-{\mdls})$ are locally free sheaves and their formations commute with any base change. Explicitly, for any morphism of $S$-schemes $u : T' \to T$, the base change morphisms are isomorphisms:
\[
    u^* \mathrm{pr}_* \cL \xrightarrow{\sim} \mathrm{pr}_* u^* \cL
\]
and
\[
    u^* \mathrm{pr}_* (\cL(-{\mdls})) \xrightarrow{\sim} \mathrm{pr}_* u^* (\cL(-{\mdls})).
\]
In particular, by \cite[Proposition~4.3]{Guignard2018}, if $\cL$ is an
invertible $\cO_C$-module of degree $d$ on each fiber of $p$, where $d$
satisfies \Cref{equation:degree-condition}, then $p_*\cL$ is a locally free
$\mathcal{O}_S$-module of rank $d-g+1$.

For further background and verification of these constructions, we refer the reader to Milne's notes on abelian Varieties (\cite{milneAV}).

\subsection{The Abel--Jacobi Morphism and Its Fibers}\label{section:AbelJacobi}


Let $U = C \setminus \mdls$ be the complement of the modulus in $C$.
The effective cartier divisors of degree $d$ which are prime to $\mdls$ are parameterized by the symmetric power $\Sym_S^d (U) = U^{(d)}$ over $S$ (See \cite{Guignard2018} Proposition 4.12, \cite{milneAV} Theorem 3.13).
For any such divisor $D \in U^{(d)}$, the associated line bundle $\Oo_C(D)$ admits a canonical trivialization along $\mdls$. 
Specifically, the canonical section $1_D$ is regular and non-vanishing on $\mdls$ because $\operatorname{supp}(D) \cap \operatorname{supp}(\mdls) = \emptyset$. 
This section restricts to a nowhere-vanishing section on the subscheme $\mathfrak{m}$, thereby determining a trivialization $\psi_D^{-1}: \Oo_C(D)|_{\mdls} \xrightarrow{\sim} \Oo_{\mdls}$.
This is done functorially in families, yielding a morphism from the symmetric power to the generalized Picard scheme (over $S$):
\begin{equation}
    \Phi_{d}: U^{(d)} \to \Pic^d_{C, \mathfrak{m}}, \quad D \mapsto [(\mathcal{O}_C(D), \psi_D)],
\end{equation}

If $\mdls=0$, $d\geq 2g-1$, and $C\to S$ admits a section, then the
Abel--Jacobi morphism $C^{(d)}\to\PicC[d]$ is a projective space bundle. In
particular, it is proper and smooth with geometrically connected fibers.

Guignard (\cite{Guignard2018} Theorem 4.14) proves that for $\mdls > 0$ and $d$ satisfying \Cref{equation:degree-condition}, the Abel-Jacobi morphism $\Phi_d$ is 
surjective smooth of relative dimension $d - \deg \mdls - g + 1$, with geometrically connected fibers. 

We now return to the standing situation $S=\Spec(k)$, which remains in
force for the rest of the chapter. The geometric fibers of $\Phi_d$ are
well understood:
\begin{theorem}\label{theorem:abel-jacobi-fibers}
Assume that $S=\Spec(k)$ and
$d \geq \max\{2g-1+\deg\mdls,\deg\mdls\}$. Let
$\bar s=\Spec\Omega\to\PicCm[d]$ be a geometric point, where $\Omega$ is
separably closed. Then the geometric fiber of the Abel--Jacobi morphism
\[
    \Phi_d:U^{(d)}\longrightarrow\PicCm[d]
\]
at $\bar s$ is isomorphic to
$$
\begin{cases}
\mathbb{A}_{\Omega}^{d - \deg \mdls - g + 1} & \text{if } \mdls > 0, \\
\mathbb{P}_{\Omega}^{d - g} & \text{if } \mdls = 0.
\end{cases}
$$
Thus every geometric fiber is an affine space when $\mdls>0$ and a
projective space when $\mdls=0$. This statement does not, by itself,
assert that $\Phi_d$ is Zariski-locally trivial.
\end{theorem}
\begin{proof}
See \cite{tendler2015geometricclassfieldtheory} Propositions 3.13-3.14, or \cite{Toth2011} Proposition 2.1.4.    

\end{proof}

\section{Takeuchi's Compactification of the Abel--Jacobi Morphism}\label{section:BlowupOfSymmetricPowerOfCurves}
We recall that our objective is to descend the local system $\cF^{[d]}$ from $U^{(d)}$ to $\PicCm[d]$ along the 
Abel-Jacobi map $\Phi_d$:

\[
\begin{tikzcd}
 \cF^{[d]} \arrow[d, purple]\\
 U^{(d)} \arrow[r, "\Phi_d"] & \PicCm[d]      
\end{tikzcd}
\]


(Here, the purple arrow emphasizes that the morphism is of sheaves on the étale site).

However, we encounter an obstruction: in the case we are considering ($\mdls > 0$), 
the fibers of $\Phi_d$ are affine spaces (of the same degree) rather than the better-behaved projective spaces. 
This suggests that a solution to this problem is to compactify the morphism to yield projective fibers. 
 
This section describes the result of the compactification constructed by \cite{takeuchi2019blow} 
by blowing up the symmetric power.

Let $\mdls = \sum_{i=1}^n k_iP_i$, where the $P_i$ are distinct closed
points of $C$, be a modulus on $C$, and let $d$ satisfy
\Cref{equation:degree-condition}. 
Takeuchi (\cite{takeuchi2019blow}) defines $Z_0=Z_0(\mathfrak{m}, d)$ as the closed subscheme of 
$C^{(d)}$ defined by the map $C^{(d - \deg \mathfrak{m})} \to C^{(d)}$ adding $\mathfrak{m}$. 
He also defines $X_{\mathfrak{m}, d}$ as the blowup of $C^{(d)}$ along $Z_0$.
Let $E_0 = E_{\mathfrak{m},d} = Z_0(\mathfrak{m}, d) \times_{C^{(d)}} X_{\mathfrak{m}, d}$ be the exceptional divisor of the blowup. It is irreducible of codimension 1, and we let $\eta_0=\eta_{\mathfrak{m}, d}$ be its generic point.


Diagrammatically:
\[\begin{tikzcd}
\overline{\{\eta_0 \}} =  E_0 \arrow[r] \arrow[d] & X_{\mdls, d} \arrow[d, "\pi"] \\
Z_0 \arrow[r, "c.i"]           & C^{(d)}          
\end{tikzcd}\] 

Incorporating $U^{(d)}$, the local system $\cF^{[d]}$ and the Abel-Jacobi map, we have:
\[
\begin{tikzcd}
    & & \cF^{[d]} \arrow[d, purple]\\
\overline{\{\eta_0 \}} =  E_0 \arrow[r] \arrow[d] & 
X_{\mdls, d} \arrow[d, "\pi"]  & \arrow[l, hook'] \arrow[dl, hook'] U^{(d)} \arrow[r, "\Phi_d"] & \PicCm[d] \\
Z_0 \arrow[r, "c.i"]           & C^{(d)}          
\end{tikzcd}
\] 

In Section 3 of \cite{takeuchi2019blow} Takeuchi constructs, for $d$ satisfying \Cref{equation:degree-condition} a compactification denoted by $\Cmod{d}$ and proves the following:

\begin{theorem}[Takeuchi]\label{theorem:takeuchi-compactification-theorem}
    The scheme $\Cmod{d}$ is an open subscheme of $X_{\mdls,d}$ containing $U^{(d)}$. 
    The morphism $\Phi_d: U^{(d)} \to \PicCm[d]$ extends to a morphism $\tilde{\Phi}_d: \Cmod{d} \to \PicCm[d]$ which makes 
    $\Cmod{d}$ a projective space bundle over $\PicCm[d]$. 
    Furthermore, the complement of $U^{(d)}$ in $\Cmod{d}$ is the scheme-theoretic intersection
    \[
        E_0 \times_{X_{\mdls,d}} \Cmod{d}
        \cong Z_0 \times_{C^{(d)}} \Cmod{d}.
    \]
\end{theorem}

     

Diagrammatically we have:
\[
\begin{tikzcd}
    &  & \cF^{[d]} \arrow[d, purple]\\
    E_0 \times_{X_{\mdls,d}} \Cmod{d} \arrow[d] \arrow[r]  & \Cmod{d} \arrow[d, hook] \arrow[rd, "\tilde{\Phi}_d"]  & U^{(d)} \arrow[l, hook'] \arrow[d, "\Phi_d"] \\
\overline{\{\eta_0 \}} =  E_0 \arrow[r] \arrow[d] & 
X_{\mdls, d} \arrow[d, "\pi"]  &  \PicCm[d] &  \\
Z_0 \arrow[r, "c.i"]           & C^{(d)}          
\end{tikzcd}
\] 

Here the boundary may equivalently be written as
$Z_0 \times_{C^{(d)}} \Cmod{d}$, since
$E_0=Z_0\times_{C^{(d)}}X_{\mdls,d}$.

The construction gives more than the description of the open complement in
\Cref{theorem:takeuchi-compactification-theorem}: it identifies the boundary
as a relative hyperplane.

\begin{cor}[The boundary is a relative hyperplane]
\label{lemma:geometric-generic-fiber-takeuchi}
The boundary
\[
    H=\Cmod{d}\setminus U^{(d)}
\]
is a smooth relative hyperplane divisor in the projective-space bundle
$\widetilde\Phi_d:\Cmod{d}\to\PicCm[d]$. Consequently, if
$\bar s=\Spec\Omega\to\PicCm[d]$ is any geometric point and
$r=d-\deg\mdls-g+1$, then there is an isomorphism of pairs
\[
    \left((\Cmod{d})_{\bar s},H_{\bar s}\right)
    \cong
    \left(\mathbb P^r_{\Omega},\mathbb P^{r-1}_{\Omega}\right),
\]
and hence $U^{(d)}_{\bar s}\cong\mathbb A^r_{\Omega}$.
\end{cor}

\begin{proof}
In Takeuchi's construction, the projective-space bundle is obtained from a
locally free sheaf $E_{\mdls}$ fitting into an exact sequence whose kernel
is $\operatorname{pr}_*\cL(-\mdls)$ and whose quotient is invertible.
The resulting closed immersion
\[
    \mathbb P\bigl(\operatorname{pr}_*\cL(-\mdls)\bigr)
    \lhook\joinrel\longrightarrow \mathbb P(E_{\mdls})
\]
is therefore a relative hyperplane. Takeuchi identifies this hyperplane
scheme-theoretically with the inverse image of $Z_0$, and identifies
$\mathbb P(E_{\mdls})$ with the blowup
\cite[Lemmas~3.2 and~3.3(1), Theorem~3.4]{takeuchi2019blow}. Pulling these
identifications back along the open immersion defining $\Cmod{d}$ gives
the asserted description of $H$. The statement on geometric fibers follows
by base change.
\end{proof}

\section{Tame Ramification Along the Full Boundary}\label{section:tame-ramification-full-boundary}

The ramification calculation of Chapter~2 concerns a modulus supported at a
single point and a symmetric power of the corresponding degree. We now
combine those local calculations over the different points of $\mdls$ and
then pass from degree $\deg\mdls$ to every sufficiently large degree $d$.

We retain the blowup notation of Chapter~2. For every nonzero effective
submodulus $\ndls\leq\mdls$, let $z_{\ndls}\in C^{(\deg\ndls)}$ be the
point representing the divisor $\ndls$, let $Z_{\ndls}$ be the
corresponding closed point, and put
\[
    X_{\ndls}=\operatorname{Bl}_{Z_{\ndls}}
        C^{(\deg\ndls)},
    \qquad
    E_{\ndls}=X_{\ndls}\times_{C^{(\deg\ndls)}}Z_{\ndls}.
\]
We write $\eta_{\ndls}$ for the generic point of $E_{\ndls}$. For the
full modulus in degree $d$, the corresponding center, blowup, exceptional
divisor, and generic point are denoted by $Z_0=Z_0(\mdls,d)$,
$X_{\mdls,d}$, $E_0=E_{\mdls,d}$, and $\eta_0=\eta_{\mdls,d}$, as in
\Cref{section:BlowupOfSymmetricPowerOfCurves}.

\begin{theorem}\label{theorem:SymmetricPowerOfSheafIsTamelyRamified}
Let $\Lambda$ be a finite commutative ring of cardinality invertible in $k$, and let
$\cF$ be an \'etale sheaf of $\Lambda$-modules, locally free of rank~$1$ on
$U$, with ramification bounded by $\mdls$. For every sufficiently large
integer $d$, the local system $\cF^{[d]}$ is tamely ramified along
\[
H=\widetilde C^{(d)}_\mdls\setminus U^{(d)}
  =E_0\times_{X_{\mdls,d}}\widetilde C^{(d)}_\mdls
  \cong Z_0\times_{C^{(d)}}\widetilde C^{(d)}_\mdls.
\]
\end{theorem}

\subsection{From Point-Supported Moduli to the Full Modulus}

The first step is to combine the single-point calculations. The following
product-compatibility result is placed here because it is used only in the
coprime-modulus reduction that immediately follows it.

\subsubsection{Compatibility under products}\label{subsection:ram_prod_analysis}


Let $X$ and $Y$ be smooth integral schemes over a field $k$, and let
$x\in X(k)$ and $y\in Y(k)$.
We denote the blowups of these schemes at the given points by 
$\pi_X: \Bl_x(X) \to X$ and $\pi_Y: \Bl_y(Y) \to Y$. Furthermore, let $\pi_{X \times Y}: \Bl_{(x,y)}(X \times_k Y) \to X \times_k Y$ 
be the blowup of the product scheme at the point $(x,y)$.
We denote by $E_X, E_Y$, and $E_{X \times Y}$ the respective exceptional divisors, and let $\eta_X, \eta_Y$, and $\eta_{X \times Y}$ be their generic points.

In this section, we establish the following result concerning the stability of ramification bounds under the external product of torsors.

\begin{prop}\label{prop:ramification-external-product}
Let $G$ be a finite abelian group. Suppose $\cG_X$ and $\cG_Y$ are $G$-torsors defined on open subsets $U_X \subset \Bl_x(X)$ and $U_Y \subset \Bl_y(Y)$ that are disjoint
from the exceptional divisors.
If the ramification of $\cG_X$ at $\eta_X$ and $\cG_Y$ at $\eta_Y$ is bounded by $r$, then the external product torsors
\[
\cG_{X \times Y} := \text{pr}_1^{-1} \cG_X \otimes \text{pr}_2^{-1} \cG_Y
\]
have ramification bounded by $r$ at the generic point $\eta_{X\times Y}$ of the exceptional divisor in the product blowup.
\end{prop}

\begin{proof}
The proposition is purely local in nature, so it suffices to consider the case where $X$ and $Y$ are affine. 
More precisely, by the smoothness of $X$ and $Y$, we may restrict our attention to open neighborhoods of $x$ and $y$ that are etale over affine spaces. 
The rest of this section treats that case.

\subsubsection*{The Affine Case}
Let $X = \mathbb{A}_k^n$ be the affine $n$-space over a field $k$, and let $0 \in X$ be the origin. 
Let $\tilde{X} = \text{Bl}_0(X)$ be the blowup of $X$ at the origin. 
Recall that $\tilde{X} \subset X \times_k \mathbb{P}^{n-1}_k$ is defined by the equations $x_i u_j = x_j u_i$, where $[u_1 : \dots : u_n]$ are the homogeneous coordinates of 
$\mathbb{P}^{n-1}_k$.
The exceptional divisor $E \subset \tilde{X}$ is the fiber over the origin, $E = \{ (0, [u_1 : \dots : u_n]) \}$, which is of codimension 1 in $\tilde{X}$. 
Let $\eta \in E$ be the generic point of $E$, and let $R = \mathcal{O}_{\tilde{X}, \eta}$ be the associated local ring. 
This ring $R$ is a discrete valuation ring (DVR) with fraction field $K = K(X) = k(x_1, \dots, x_n)$.

On the affine chart $U_1$ where $u_1 \neq 0$, we have $x_i = \frac{u_i}{u_1} x_1$. 
The coordinate ring is:$$\mathcal{O}_{\tilde{X}}(U_1) = k\left[x_1, \frac{u_2}{u_1}, \dots, \frac{u_n}{u_1}\right]$$
In this chart, the generic point $\eta$ corresponds to the prime ideal $\mathfrak{p}_1 = (x_1)$. Thus, the local ring is $R = k[x_1, \frac{u_2}{u_1}, \dots, \frac{u_n}{u_1}]_{(x_1)}$. 
The residue field is $\kappa(\eta) = k(\frac{u_2}{u_1}, \dots, \frac{u_n}{u_1})$, and the completion of $R$ with respect to its maximal ideal is:
$$\hat{R} = \kappa(\eta)[[x_1]]$$
In this local ring, $x_1$ is a uniformizer. Note that any $x_i$ (for $i > 1$) can also serve as a uniformizer, 
as $x_i = (\frac{u_i}{u_1}) x_1$ and $\frac{u_i}{u_1}$ is a unit in $R$.

For any monomial $M = x_1^{a_1} \dots x_n^{a_n} \in K$, we can write:
$$M = x_1^{\sum a_i} \left( \frac{u_2}{u_1} \right)^{a_2} \dots \left( \frac{u_n}{u_1} \right)^{a_n}$$
Since the term in parentheses is a unit in $R$, the valuation $\nu_E$ associated with $E$ satisfies:
$$\nu_E(M) = \sum a_i = \deg M$$
Consequently, for any polynomial $f = f_d + f_{d+1} + \dots + f_l$, where $f_i$ is the homogeneous part of degree $i$, we have $\nu_E(f) = d$ (the order of vanishing at the origin).

\subsubsection*{The Product Case}
Now, let $X = \mathbb{A}^n$ and $Y = \mathbb{A}^m$ with origins $x=0$ and $y=0$. 
As before, $\text{Bl}_0(X) \subset X \times \mathbb{P}^{n-1}$ and $\text{Bl}_0(Y) \subset Y \times \mathbb{P}^{m-1}$ have exceptional divisors $E_X$ and $E_Y$ respectively.
Consider the product $X \times_k Y \cong \mathbb{A}_k^{n+m}$. 
The blowup of the product at the origin $(0,0)$, denoted $\text{Bl}_{(0,0)}(X \times_k Y)$, is a subscheme of $(X \times Y) \times \mathbb{P}^{n+m-1}$ defined by:

$$\begin{cases} 
x_i w_j = x_j w_i & 1 \le i, j \le n \\
y_k w_{n+l} = y_l w_{n+k} & 1 \le k, l \le m \\
x_i w_{n+j} = y_j w_i & 1 \le i \le n, \,\, 1 \le j \le m 
\end{cases}$$
where $[w_1 : \dots : w_{n+m}]$ are the homogeneous coordinates of $\mathbb{P}^{n+m-1}$. The exceptional divisor $E_{X \times Y}$ is isomorphic to $\mathbb{P}^{n+m-1}$.

\subsubsection*{Comparison of Blowups}

Both $\text{Bl}_0(X) \times_k \text{Bl}_0(Y)$ and
$\text{Bl}_{(0,0)}(X \times_k Y)$ are birational to $X \times_k Y$; their
common dense open is
\[
    (X\setminus\{0\})\times_k(Y\setminus\{0\}).
\]
Put $B=\text{Bl}_{(0,0)}(X\times_kY)$, and let
\[
    q\colon B\longrightarrow\mathbb P_k^{n+m-1}
\]
be the second projection in the above realization of $B$. Define
\[
    \Omega_X=\bigcup_{i=1}^{n}D_+(w_i),
    \qquad
    \Omega_Y=\bigcup_{j=1}^{m}D_+(w_{n+j}),
\]
and
\[
    \widetilde U
      :=q^{-1}(\Omega_X\cap\Omega_Y)
      =\bigcup_{\substack{1\leq i\leq n\\1\leq j\leq m}}
        q^{-1}\!\left(D_+(w_iw_{n+j})\right).
\]
There is a morphism
\[
\begin{aligned}
f\colon\widetilde U&\longrightarrow
   \text{Bl}_0(X)\times_k\text{Bl}_0(Y),\\
((x,y),[w_1:\cdots:w_{n+m}])&\longmapsto
\bigl((x,[w_1:\cdots:w_n]),
      (y,[w_{n+1}:\cdots:w_{n+m}])\bigr).
\end{aligned}
\]
Thus we have a diagram
\[
\begin{tikzcd}
& \widetilde U \arrow[dl,"f"'] \arrow[dr,hook,"j"] & \\
\text{Bl}_0(X)\times_k\text{Bl}_0(Y)
&& B,
\end{tikzcd}
\]
where $j$ is the open immersion. The restriction of $q$ induces an
isomorphism $E_{X\times Y}\xrightarrow{\sim}\mathbb P_k^{n+m-1}$. Since neither
$V_+(w_1,\ldots,w_n)$ nor $V_+(w_{n+1},\ldots,w_{n+m})$ contains the generic
point of $\mathbb P_k^{n+m-1}$, we have
$\eta_{X\times Y}\in\widetilde U$.

\subsubsection*{Extensions of DVRs}
Let
\[
S=\mathcal O_{B,\eta_{X\times Y}},\qquad
R_X=\mathcal O_{\text{Bl}_0(X),\eta_X},\qquad
R_Y=\mathcal O_{\text{Bl}_0(Y),\eta_Y}.
\]
The morphisms $\operatorname{pr}_X\circ f$ and
$\operatorname{pr}_Y\circ f$ send $\eta_{X\times Y}$ to $\eta_X$ and
$\eta_Y$, respectively, and therefore induce local homomorphisms of DVRs
\[
    R_X\longrightarrow S,\qquad R_Y\longrightarrow S.
\]
The point $\eta_{X\times Y}$ belongs to
$q^{-1}(D_+(w_1w_{n+1}))$. On this chart,
\[
\begin{aligned}
S
&=k\left[
x_1,\frac{x_2}{x_1},\ldots,\frac{x_n}{x_1},
\frac{y_1}{x_1},\ldots,\frac{y_m}{x_1}
\right]_{(x_1)}\\
&=k\left[
y_1,\frac{x_1}{y_1},\ldots,\frac{x_n}{y_1},
\frac{y_2}{y_1},\ldots,\frac{y_m}{y_1}
\right]_{(y_1)}.
\end{aligned}
\]
Consequently, $x_1$ is a uniformizer in both $R_X$ and $S$, while $y_1$ is
a uniformizer in both $R_Y$ and $S$. Both extensions have ramification index
one. Their residue-field extensions are
\[
k\left(\frac{x_2}{x_1},\ldots,\frac{x_n}{x_1}\right)
\subset
k\left(\frac{x_2}{x_1},\ldots,\frac{x_n}{x_1},
        \frac{y_1}{x_1},\ldots,\frac{y_m}{x_1}\right)
\]
and
\[
k\left(\frac{y_2}{y_1},\ldots,\frac{y_m}{y_1}\right)
\subset
k\left(\frac{x_1}{y_1},\ldots,\frac{x_n}{y_1},
        \frac{y_2}{y_1},\ldots,\frac{y_m}{y_1}\right),
\]
respectively. In particular, both are purely transcendental and hence
separable.

\subsubsection*{Ramification of $G$-Torsors}
The extensions $R_X\to S$ and $R_Y\to S$ are weakly unramified and have
separable residue-field extensions. Hence the pullbacks of $\cG_X$ and
$\cG_Y$ to $\operatorname{Frac}(S)$ both have ramification bounded by $r$.
Their contracted product is the restriction of $\cG_{X\times Y}$ to
$\operatorname{Frac}(S)$, so
\Cref{lemma:bounding_ramification_of_contracted_product} gives the same bound
for $\cG_{X\times Y}$ at $\eta_{X\times Y}$.
\end{proof}

\begin{lemma}\label{lemma:moduli_reduction}
     Let $\ndls_1, \ndls_2 \subset \mdls$ be two coprime nonzero
     submoduli. Assume that $\cP^{[\deg \ndls_1]}$ and
     $\cP^{[\deg \ndls_2]}$ are tamely ramified at $\eta_{\ndls_1}$ and
     $\eta_{\ndls_2}$, respectively. Then
     $\cP^{[\deg \ndls_1 + \deg \ndls_2]}$ is tamely ramified at
     $\eta_{\ndls_1 + \ndls_2}$.
\end{lemma}

\begin{proof}
Apply \Cref{prop:ramification-external-product} to the two torsors and then
use the compatibility of symmetric powers with addition of divisors from
\Cref{prop:symmetric_powers_on_curves}.
\end{proof}

\subsection{From the Minimal Degree to Large Degrees}

We now propagate the result in degree $\deg\mdls$ to degree $d$ by adding
an arbitrary effective divisor of degree $d-\deg\mdls$.

The following addition lemma is adapted from \cite[Lemma~4.1]{takeuchi2019blow}.

\begin{lemma}\label{lemma:takeuchi_lemma}
     Let $C$ be a projective, smooth, and geometrically connected curve over a perfect field $k$. 
     Let $\mathfrak{m} = \sum_{i=1}^{r} k_i P_i$ be an effective divisor where $P_1, \dots, P_r$ are distinct closed points. 
     Let $U = C \setminus \mathfrak{m}$ and let $d \geq \deg \mathfrak{m}$.

     Suppose $\mathfrak{n}_1, \dots, \mathfrak{n}_l$ are pairwise coprime submoduli of $\mathfrak{m}$ such that $\mathfrak{m} = \sum_{j=1}^l \mathfrak{n}_j$. 
     Consider the summation morphism:
     \[
     \pi : C^{(\deg \mathfrak{n}_1)} \times_k \dots \times_k C^{(\deg \mathfrak{n}_l)} \times_k C^{(d - \deg \mathfrak{m})} \longrightarrow C^{(d)}
     \]
     defined by $(D_1, \dots, D_l, D_{extra}) \mapsto \sum_{j=1}^l D_j + D_{extra}$.

     Then $\pi$ is étale at the generic point of the closed subvariety 
     \[
     V = \{\mathfrak{n}_1\} \times_k \dots \times_k \{\mathfrak{n}_l\} \times_k C^{(d - \deg \mathfrak{m})}
     \]
     inside the domain $C^{(\deg \mathfrak{n}_1)} \times_k \dots \times_k C^{(\deg \mathfrak{n}_l)} \times_k C^{(d - \deg \mathfrak{m})}$.

\begin{proof}
     We may assume that $k$ is algebraically closed (hence $\deg P_i = 1$ for all $i$).
     The morphism $\pi$ is finite: it is induced by the finite quotient morphism
     \[
     C^d/\bigl(\mathfrak S_{\deg\ndls_1}\times\cdots\times
       \mathfrak S_{\deg\ndls_l}\times
       \mathfrak S_{d-\deg\mdls}\bigr)
       \longrightarrow C^d/\mathfrak S_d.
     \]
     Its source and target are smooth $k$-schemes of pure dimension $d$.
     Hence $\pi$ is flat by miracle flatness, and therefore finite flat.
     It is enough to show that there exists a closed point $Q$ of $\ndls_1 + \dots \ndls_l + C^{(d - \deg \mathfrak{m})} \subset C^{(d)}$ 
     over which there are $\deg \pi$ points on $C^{(\deg\ndls_1)} \times_k \dots \times_k C^{(\deg\ndls_l)} \times_k C^{(d - \deg \mathfrak{m})}$. (Because it will be unramified 
     at this point and thus also at the generic point of $V$.)
     Choose $Q$ as a point corresponding to a divisor
     $\ndls_1 + \dots + \ndls_l + P_{r+1} + \dots + P_{r+d - \deg \mathfrak{m}}$, where $P_1, \dots, P_{r+d - \deg \mathfrak{m}}$ are distinct points of $U(k)$.
\end{proof}     
\end{lemma}

\begin{cor}\label{cor:etale_at_generic_point}
     The morphism $C^{(\deg \mdls)} \times_k C^{(d - \deg \mdls)} \xrightarrow[]{\pi} C^{(d)}$
     is finite flat everywhere, and étale at the generic point of the closed subvariety
     $Z_\mdls \times_k C^{(d - \deg \mdls)} \subset C^{(\deg \mdls)} \times_k C^{(d - \deg \mdls)}$.     
\end{cor}
\begin{proof}
This is the case $l=1$ and $\mathfrak n_1=\mathfrak m$ of \Cref{lemma:takeuchi_lemma}.
\end{proof}

Following from this, we look at the following diagram, coming from the flat base change

$C^{(d - \deg \mdls)} \to \Spec(k)$ (\Cref{prop:symmetric_power_flat_base_change}) of \Cref{equation:blowup_diagram_ndls}:

\[\begin{tikzcd}
E_{\mdls} \times_k C^{(d-\deg \mdls)} \arrow[r] \arrow[d] & X_{\mdls} \times_k C^{(d-\deg \mdls)} \arrow[d] & \cF^{[\deg \mdls]} \times_k C^{(d-\deg \mdls)} \arrow[d, color=purple] \\
Z_\mdls \times_k C^{(d-\deg \mdls)} \arrow[r] & C^{(\deg \mdls)} \times_k C^{(d-\deg \mdls)} & U^{(\deg \mdls)} \times_k C^{(d-\deg \mdls)} \arrow[lu]
\end{tikzcd}\]

Note that $U^{(\deg \mdls)} \times_k C^{(d-\deg \mdls)}$ is dense open subscheme of $X_{\mdls} \times_k C^{(d-\deg \mdls)}$,
And $E_{\mdls} \times_k C^{(d-\deg \mdls)}$ is a prime divisor of $X_{\mdls} \times_k C^{(d-\deg \mdls)}$.
Hence it is a well-defined question according to \Cref{definition:tamely_ramified_in_codim_1} to ask whether
$\cF^{[\deg \mdls]} \times_k C^{(d-\deg \mdls)}$ is tamely ramified at the generic point $\theta$ of
$E_{\mdls} \times_k C^{(d-\deg \mdls)}$.

\begin{lemma}\label{lemma:tame_ramification_persists}
     If $\cF^{[\deg \mdls]}$ is tamely ramified at $\eta_\mdls$, then
     $\cF^{[\deg \mdls]} \times_k C^{(d-\deg \mdls)}$ is tamely ramified at the generic point $\theta$ of
     $E_{\mdls} \times_k C^{(d-\deg \mdls)}$.
\end{lemma}
\begin{proof}
This is the stability of tame ramification in codimension one under base
change (\stackstag{0EYD}), applied to the projection
\[
X_{\mdls}\times_k C^{(d-\deg\mdls)}\longrightarrow X_{\mdls}.
\]
\end{proof}


By \Cref{lemma:tame_ramification_persists}, replacing $C^{(d - \deg \mdls)}$ with the dense open subscheme $U^{(d - \deg \mdls)} \subset C^{(d - \deg \mdls)}$,
we get that the $G$-torsor $p_1^{-1} \cP^{[\deg \mdls]}$, where
$\cP=\operatorname{Fr}(\cF)$,
is tamely ramified at $\theta$ the generic point of
$E_{\mdls} \times_k U^{(d-\deg \mdls)}$, viewed as a prime divisor of
$X_{\mdls}\times_k U^{(d-\deg\mdls)}$; the torsor itself is defined on the
dense open $U^{(\deg\mdls)}\times_k U^{(d-\deg\mdls)}$,
where $p_1: U^{(\deg \mdls)} \times_k U^{(d - \deg \mdls)} \to U^{(\deg \mdls)}$ is the projection to the first factor.

Looking at the second projection $p_2:X_{\mdls} \times_k U^{(d - \deg \mdls)} \to U^{(d - \deg \mdls)}$, and the fact that 
$\cP^{[d - \deg \mdls]}$ is \'etale on $U^{(d - \deg \mdls)}$ we get that 
$p_2^{-1} \cP^{[d - \deg \mdls]}= X_{\mdls} \times_k \cP^{[d - \deg \mdls]}$ is \'etale on $X_{\mdls} \times_k U^{(d - \deg \mdls)}$.
Hence, its restriction to $U^{(\deg \mdls)} \times_k U^{(d - \deg \mdls)}$ is unramified at $\theta$.

Thus, by \Cref{lemma:tame-times-unramified}, we conclude that $\boxP{\cP^{[\deg \mdls]}}{\cP^{[d - \deg \mdls]}} = p_1^{-1} \cP^{[\deg \mdls]} \otimes^G p_2^{-1} \cP^{[d - \deg \mdls]}$
     is tamely ramified at $\theta$ the generic point of $E_{\mdls} \times_k U^{(d-\deg \mdls)}$.

\begin{lemma}\label{lemma:tame-times-unramified}
     Let $X \to \Spec k$ be a scheme over a field $k$, and let $\cP_1, \cP_2$ be two $G$-torsors on $U_{et}$.
     Let $\xi$ be the generic point of a prime divisor $D \subset X$.
     If $\cP_1$ is tamely ramified at $\xi$, and $\cP_2$ is unramified at $\xi$,
     then the product $\cP_1 \otimes^G \cP_2$ is tamely ramified at $\xi$.
\end{lemma}
\begin{proof}
     This follows from \Cref{lemma:bounding_ramification_of_contracted_product}.
\end{proof}

Combining this with \Cref{cor:etale_at_generic_point} we get

\begin{cor}\label{cor:tame-ramification-symmetric-product}
     Let $\mdls$ be a modulus as above, and let $\eta_\mdls$ be the generic point of $E_\mdls$.
     Let $\cP$ be a $G$-torsor on $U_{et}$ with ramification bounded by $\mdls$.
     Assume $\cP^{[\deg \mdls]}$ is tamely ramified at $\eta_\mdls$.
     Then $\boxP{\cP^{[\deg \mdls]}}{\cP^{[d - \deg \mdls]}}$ is tamely ramified at the generic point $\theta$ of 
     $E_{\mdls} \times_k C^{(d-\deg \mdls)} \subset C^{(\deg \mdls)} \times_k C^{(d - \deg \mdls)}$,
     and $\cP^{[d]}$ is tamely ramified at the generic point $\eta_0=\eta_{\mdls,d}$ of $E_0=E_{\mdls,d}$
\end{cor}
\begin{proof} 
The first assertion, that $\mathcal{P}^{[\deg \mathfrak{m}]} \boxtimes \mathcal{P}^{[d - \deg \mathfrak{m}]}$ is tamely ramified at $\theta$, follows from the preceding discussion. Thus, it remains to show that $\mathcal{P}^{[d]}$ is tamely ramified at $\eta_0$.

Consider the blowup diagram defining $X_{\mathfrak{m}, d}$:
\[
\begin{tikzcd}
\overline{\{\eta_0 \}} =  E_0 \arrow[r] \arrow[d] & X_{\mdls, d} \arrow[d, "\pi"] \\
Z_0 \arrow[r, "c.i"]           & C^{(d)}          
\end{tikzcd}
\] 
By performing a base change along the flat addition map $+: C^{(\deg \mathfrak{m})} \times_k U^{(d-\deg \mathfrak{m})} \to C^{(d)}$, we obtain the following commutative diagram:

\[
\begin{tikzcd}
&  \left(C^{(\deg \mdls)} \times_k U^{(d-\deg \mdls)} \right) \times_{C^{(d)}} E_0 \arrow[r] \arrow[d] & \overline{\{\eta_0 \}} =  E_0 \arrow[d] \\
& \left(C^{(\deg \mdls)} \times_k U^{(d-\deg \mdls)} \right) \times_{C^{(d)}} X_{\mdls, d} \arrow[d] \arrow[r] & X_{\mdls, d} \arrow[d, "\pi"] \\
\mdls \times_k U^{(d-\deg \mdls)} \arrow[r, "c.i"] & C^{(\deg \mdls)} \times_k U^{(d-\deg \mdls)} \arrow[r, "+"]  & C^{(d)}          
\end{tikzcd}
\] 
By \Cref{theorem:blowup_flat_base_change}, blowups commute with flat base
change. Since the inverse image of the center $Z_0$ under the map $+$ is
$\mathfrak{m} \times_k U^{(d-\deg \mathfrak{m})}$, the scheme
$\left(C^{(\deg \mathfrak{m})} \times_k U^{(d-\deg \mathfrak{m})} \right)
\times_{C^{(d)}} X_{\mathfrak{m}, d}$ is the blowup of
$C^{(\deg \mathfrak{m})} \times_k U^{(d-\deg \mathfrak{m})}$ along
$\mathfrak{m} \times_k U^{(d-\deg \mathfrak{m})}$.
This, in turn, is isomorphic to the base change of the blowup $X_{\mathfrak{m}}$ (of $C^{(\deg \mathfrak{m})}$ along $\mathfrak{m}$) via the (flat) projection $C^{(\deg \mathfrak{m})} \times_k U^{(d - \deg \mathfrak{m})} \to C^{(\deg \mathfrak{m})}$.

Assembling these facts, we obtain the following Cartesian square:
\[
\begin{tikzcd}
&  E_{\mdls}\times U^{(d - \deg \mdls)}  \arrow[r, "\tilde{+}"] \arrow[d] & \overline{\{\eta_0 \}} =  E_0 \arrow[d] \\
& X_{\mdls} \times_k U^{(d - \deg \mdls)} \arrow[d] \arrow[r, "\tilde{+}"] & X_{\mdls, d} \arrow[d, "\pi"] \\
\mdls \times_k U^{(d-\deg \mdls)} \arrow[r, "c.i"] & C^{(\deg \mdls)} \times_k U^{(d-\deg \mdls)} \arrow[r, "+"]  & C^{(d)}          
\end{tikzcd}
\] 

The point $\theta$ defined in the Corollary is the generic point of $E_{\mathfrak{m}} \times_k U^{(d-\deg \mathfrak{m})}$.
 Let $\eta$ be the generic point of $\mathfrak{m} \times_k U^{(d-\deg \mathfrak{m})}$. 
 By \Cref{cor:etale_at_generic_point}, the map $+$ is étale at $\eta$. 
Consequently, the lifted map $\tilde{+}$ is étale at $\theta$. 
Given the isomorphism $(+^{-1})(\cP^{[d]}) \cong \cP^{[\deg \mdls]} \boxtimes \cP^{[d - \deg \mdls]}$ from \Cref{prop:symmetric_powers_on_curves},
the tame ramification of the box product at $\theta$ descends to the tame ramification of $\mathcal{P}^{[d]}$ at $\eta_0$ by applying \Cref{lemma:descend_ramification_along_etale}.

\end{proof}


\subsection{Conclusion of the Tameness Argument}

\begin{proof}[Proof of \Cref{theorem:SymmetricPowerOfSheafIsTamelyRamified}]
Let $\cP=\operatorname{Fr}(\cF)$ be the $\Lambda^\times$-torsor of frames.
For each single-point submodulus $\ndls=k_P P$ of $\mdls$,
\Cref{theorem:SymmetricPowerOfSheavesIsTamelyRamifiedReduction} shows that
$\cP^{[\deg\ndls]}$ is tamely ramified at $\eta_\ndls$. Repeated application
of \Cref{lemma:moduli_reduction} shows that
$\cP^{[\deg\mdls]}$ is tamely ramified at $\eta_\mdls$. Therefore
\Cref{cor:tame-ramification-symmetric-product} implies that $\cP^{[d]}$ is
tamely ramified at the generic point $\eta_0$ of $E_0$ for every
sufficiently large $d$. The same inertia character governs $\cF^{[d]}$, so
the same conclusion holds for $\cF^{[d]}$. Since $H$ is the nonempty open
part of $E_0$ described in
\Cref{lemma:geometric-generic-fiber-takeuchi}, its generic point is
$\eta_0$. Thus $\cF^{[d]}$ is tamely ramified along $H$.
\end{proof}

\section{Extension Across the Boundary}\label{section:extension-across-boundary}

We now use the geometry of the Abel--Jacobi fibers to upgrade the tameness
proved above to extension across the boundary. On the geometric generic
fiber, a Tendler-style prime-to-$p$ power argument forces the boundary
inertia representation to be trivial.

\begin{lemma}\label{lemma:ExtensionOfTamelyRamifiedSheaf}
    Let $\Lambda$ be a finite commutative ring whose cardinality is
    invertible in $k$, and let $\cF^{[d]}$ be a locally constant sheaf of
    free rank-one $\Lambda$-modules on $U^{(d)}$. If $\cF^{[d]}$ is tamely
    ramified at the generic point of
    \[
        H=\widetilde C^{(d)}_\mdls\setminus U^{(d)},
    \]
    then it extends uniquely to a locally constant sheaf
    $\widetilde{\cF}^{[d]}$ on $\widetilde C^{(d)}_\mdls$.
\end{lemma}

\begin{proof}
Put $r=d-\deg\mdls-g+1$, let $\xi$ be the generic point of $\PicCm[d]$,
let $\Omega=\kappa(\xi)^{\mathrm{sep}}$, and write
$\bar\xi=\Spec\Omega$. By
\Cref{lemma:geometric-generic-fiber-takeuchi}, the compactified geometric
generic Abel--Jacobi fiber and its boundary are
\[
    \left(
    \widetilde C^{(d)}_{\mdls,\bar\xi},
    H_{\bar\xi}
    \right)
    \cong
    \left(
    \mathbb P^r_{\Omega},
    \mathbb P^{r-1}_{\Omega}
    \right),
\]
and hence the geometric generic open fiber is
\[
    U^{(d)}_{\bar\xi}\cong\mathbb A^r_{\Omega}.
\]

Let
\[
    \chi\colon
    \pi_1\bigl(\mathbb A^r_{\Omega}\bigr)
       \longrightarrow \Lambda^\times
\]
be the character corresponding to the restriction of $\cF^{[d]}$ to this
fiber.

Let $\eta$ be the generic point of $H$, and let $\bar\eta$ be the generic
point of $H_{\bar\xi}$. Because $H$ is a relative hyperplane divisor, the
passage from $\eta$ to $\bar\eta$ preserves a local equation of the
boundary and induces a separable extension of residue fields. The
corresponding extension of henselian discrete valuation rings is therefore
unramified, and it identifies the inertia group at $\eta$ with the inertia
group
\[
    I_\infty:=I_{\bar\eta}
\]
at the hyperplane at infinity. In particular, tameness and triviality of
the finite inertia representation may both be checked after this base
change.

If $\operatorname{char}(k)=0$, then
\[
    \pi_1\bigl(\mathbb A^r_{\Omega}\bigr)=1,
\]
so $\chi$ is trivial
\cite[Example~4.14]{tendler2015geometricclassfieldtheory}.

Suppose now that $\operatorname{char}(k)=p>0$. Since $\cF^{[d]}$ is tame
along $H$,
the finite group $\chi(I_\infty)$ has order prime to $p$. Let $N$ be its
exponent. Then $(N,p)=1$, and the character $\chi^N$ is trivial on
$I_\infty$.

Consequently, the local system corresponding to $\chi^N$ is unramified at
the generic point of $\mathbb P^{r-1}_{\Omega}$. By purity, it extends to
$\mathbb P^r_{\Omega}$. Since
\[
    \pi_1\bigl(\mathbb P^r_{\Omega}\bigr)=1,
\]
this extension is constant, and therefore
\[
    \chi^N=1.
\]

On the other hand, $\chi$ factors through
\[
    \pi_1\bigl(\mathbb A^r_{\Omega}\bigr)^{\mathrm{ab}},
\]
which is a pro-$p$ group
\cite[Example~4.15]{tendler2015geometricclassfieldtheory}. Hence the finite
group $\operatorname{im}(\chi)$ is a $p$-group. Since its exponent divides
$N$, and $N$ is prime to $p$, it follows that
\[
    \operatorname{im}(\chi)=1.
\]
Thus $\chi$ is trivial.

In either characteristic, the restriction of $\chi$ to $I_\infty$ is
trivial. By the inertia comparison above, the inertia group at $\eta$ acts
trivially on $\cF^{[d]}$.

Therefore $\cF^{[d]}$ is unramified at the generic point of $H$. Since
$\widetilde C^{(d)}_\mdls$ is regular and its complement of $U^{(d)}$ is
the single smooth relative hyperplane divisor $H$, purity of the branch
locus (\stackstag{0BMB}) gives an extension of $\cF^{[d]}$ to
$\widetilde C^{(d)}_\mdls$. The extension is unique because
\[
    \pi_1\bigl(U^{(d)}\bigr)
       \longrightarrow
    \pi_1\bigl(\widetilde C^{(d)}_\mdls\bigr)
\]
is surjective.
\end{proof}

\begin{rem*}
The argument does not assert that tame ramification is generally
unramified. Tameness is used only to choose $N$ prime to $p$. The conclusion
that the inertia character is trivial then follows from the global geometry
of the geometric generic fiber and from the fact that the abelian
fundamental group of affine space is pro-$p$.
\end{rem*}

\section{Descent and Reduced Geometric Class Field Theory}\label{section:gcft-descent}

We finish by recording the fundamental-group input and applying it to the
ordinary and generalized Abel--Jacobi morphisms.

\subsection{The Fundamental-Group Input}

\begin{prop}[\stackstag{0C0J}]\label{prop:fundamental-group-exact-sequence}
Let $f\colon X\to S$ be a flat proper morphism of finite presentation whose
geometric fibers are connected and reduced. Assume that $S$ is connected,
and let $\bar s$ be a geometric point of $S$. Then there is an exact
sequence
\[
\pi_1(X_{\bar s})\longrightarrow\pi_1(X)\longrightarrow\pi_1(S)
\longrightarrow 1.
\]
\end{prop}

\begin{cor}\label{cor:etale_fundamental_group_smooth_proper}
Let $f\colon X\to S$ be a proper smooth morphism of finite presentation
whose geometric fibers are connected. Then the sequence of
\Cref{prop:fundamental-group-exact-sequence} is exact.
\end{cor}
\begin{proof}
A smooth geometric fiber is reduced, so this is an immediate application of
\Cref{prop:fundamental-group-exact-sequence}.
\end{proof}

\begin{cor}\label{cor:projective-bundle-fundamental-group}
Let $f\colon X\to S$ be a projective space bundle over a connected scheme.
Then pullback induces an isomorphism
\[
\pi_1^{\mathrm{et}}(X)\xrightarrow{\sim}\pi_1^{\mathrm{et}}(S).
\]
\end{cor}
\begin{proof}
The geometric fibers are projective spaces and therefore have trivial
\'etale fundamental group. The assertion follows from
\Cref{cor:etale_fundamental_group_smooth_proper}.
\end{proof}

\subsection{The Descent Argument}

We work over $S=\Spec k$, where $k$ is perfect. It is convenient first to
prove the following torsor formulation.

\begin{theorem}[Torsor formulation]\label{thm:GCFT_torsors}
Let $G=\Lambda^\times$, and let $\cP$ be a $G$-torsor on $U$ whose
ramification is bounded by $\mdls$. For every sufficiently large integer
$d$, there exists a $G$-torsor $\cQ_d$ on $\PicCm[d]$, unique up to
isomorphism, such that
\[
\Phi_d^*\cQ_d\cong\cP^{[d]}.
\]
\end{theorem}

\begin{proof}
We divide the proof into two cases.

\textbf{Case 1: $\mdls=0$.}
For $d$ sufficiently large,
the usual Abel--Jacobi morphism
\[
\Phi_d\colon C^{(d)}\longrightarrow\PicC[d]
\]
is proper and smooth with geometrically connected fibers, and
\Cref{theorem:abel-jacobi-fibers} identifies every geometric fiber with
$\mathbb P^{d-g}$. Applying
\Cref{prop:fundamental-group-exact-sequence}, and using
$\pi_1^{\mathrm{et}}(\mathbb P^{d-g})=1$, gives
\[
\pi_1^{\mathrm{et}}\bigl(C^{(d)}\bigr)
 \xrightarrow{\sim}
\pi_1^{\mathrm{et}}\bigl(\PicC[d]\bigr).
\]
The $G$-torsor $\cP^{[d]}$ therefore descends uniquely to a $G$-torsor on
$\PicC[d]$.

\textbf{Case 2: $\mdls>0$.}
Let $\cF=\cP\times^G\Lambda$ be the rank-one local system associated with
$\cP$; its symmetric power has frame torsor $\cP^{[d]}$. For $d$
sufficiently large,
\Cref{theorem:takeuchi-compactification-theorem} extends $\Phi_d$ to a
projective space bundle
\[
\widetilde\Phi_d\colon
\widetilde C^{(d)}_\mdls\longrightarrow\PicCm[d].
\]
By \Cref{theorem:SymmetricPowerOfSheafIsTamelyRamified},
$\cF^{[d]}$ is tamely ramified along
$H=\widetilde C^{(d)}_\mdls\setminus U^{(d)}$. Hence
\Cref{lemma:ExtensionOfTamelyRamifiedSheaf} gives a unique extension
$\widetilde{\cF}^{[d]}$ to $\widetilde C^{(d)}_\mdls$.

By \Cref{cor:projective-bundle-fundamental-group},
\[
\pi_1^{\mathrm{et}}\bigl(\widetilde C^{(d)}_\mdls\bigr)
 \xrightarrow{\sim}
\pi_1^{\mathrm{et}}\bigl(\PicCm[d]\bigr).
\]
Thus $\widetilde{\cF}^{[d]}$ descends uniquely to a rank-one local system
$\cG_d$ on $\PicCm[d]$. Restricting to $U^{(d)}$ gives
$\Phi_d^*\cG_d\cong\cF^{[d]}$. Taking frame torsors yields
\[
\Phi_d^*\operatorname{Fr}(\cG_d)\cong\cP^{[d]},
\]
so $\cQ_d=\operatorname{Fr}(\cG_d)$ has the required property. Uniqueness
follows from the fundamental-group isomorphism and from the surjection
\[
\pi_1^{\mathrm{et}}(U^{(d)})\twoheadrightarrow
\pi_1^{\mathrm{et}}(\widetilde C^{(d)}_\mdls)
\]
induced by the dense open immersion into the normal scheme
$\widetilde C^{(d)}_\mdls$.
\end{proof}

\begin{proof}[Proof of \Cref{thm:GCFT_reduced}]
Let $\cF$ be as in \Cref{thm:GCFT_reduced}, and set
$\cP=\operatorname{Fr}(\cF)$. Apply \Cref{thm:GCFT_torsors} to obtain
$\cQ_d$ on $\PicCm[d]$, and let
\[
\cG_d=\cQ_d\times^{\Lambda^\times}\Lambda.
\]
The associated-module construction commutes with pullback and carries
$\cP^{[d]}$ to $\cF^{[d]}$, so
$\Phi_d^*\cG_d\cong\cF^{[d]}$. If $\cG'_d$ is another such local system,
then $\operatorname{Fr}(\cG'_d)$ has the defining property of $\cQ_d$.
The uniqueness of $\cQ_d$, followed by the associated-module construction,
gives $\cG'_d\cong\cG_d$.
\end{proof}
